% !TEX program = xelatex
\documentclass[11pt,a4paper]{ctexart}
\usepackage[margin=2.2cm]{geometry}
\usepackage{amsmath,amssymb}
\usepackage{xcolor}
\usepackage{enumitem}
\usepackage{booktabs}
\usepackage{longtable}
\usepackage{array}
\usepackage{xurl}
\usepackage{hyperref}
\hypersetup{colorlinks=true,linkcolor=blue!50!black,urlcolor=blue!50!black}

\setlength{\parskip}{3pt}
\newcounter{finding}
\newenvironment{finding}[1]{%
  \refstepcounter{finding}%
  \subsection*{条目 \thefinding\quad #1}%
  \addcontentsline{toc}{subsection}{条目 \thefinding\quad #1}%
}{\par\medskip\noindent\rule{\linewidth}{0.3pt}}
\newcommand{\loc}[1]{\par\noindent\textbf{【原文位置】}#1}
\newcommand{\orig}[1]{\par\noindent\textbf{【原文内容】}\begin{quote}\small\itshape #1\end{quote}}
\newcommand{\ptype}[1]{\par\noindent\textbf{【问题类型】}#1}
\newcommand{\basis}[1]{\par\noindent\textbf{【具体依据及对理解或论证的影响】}#1}
\newcommand{\conf}[1]{\par\noindent\textbf{【判断把握】}#1}
\newcommand{\normal}[1]{\par\noindent\textbf{【是否可能属于正常学术表达及原因】}#1}
\newcommand{\advice}[1]{\par\noindent\textbf{【建议】}#1}
\newcommand{\revise}[1]{\par\noindent\textbf{【建议修改文本】}\begin{quote}\small #1\end{quote}}
\newcommand{\authfill}[1]{\textcolor{red!70!black}{[待作者补充：#1]}}
\newcommand{\verify}{\textcolor{orange!80!black}{\textbf{待核实}}}
% 建议修改文本中保留原稿的 \citep / \ref 命令，按原样显示，便于直接粘回 main.tex
\newcommand{\citep}[1]{\texttt{\textbackslash citep\{}\path{#1}\texttt{\}}}
\newcommand{\bibkey}[1]{\path{#1}}
\setlength{\emergencystretch}{3em}
\AtBeginDocument{\renewcommand{\ref}[1]{\texttt{\textbackslash ref\{#1\}}}}

\title{M4Fed 稿件写作特征检查报告\\[2pt]\large （模板化、空泛、重复、证据边界与一致性）}
\author{审查对象：\texttt{latex/main.tex}、\texttt{latex/supplementary.tex}}
\date{2026年11月6日}

\begin{document}
\maketitle

\begin{abstract}
本报告检查稿件 \emph{A Multi-Source Federated Learning Framework with Medical Foundation Models for Multi-Modal Medical Data on Brain Age Estimation} 中模板化、空泛、机械重复或表达不自然的句段，并核对论证与证据的对应关系和文内一致性。报告只分析写作特征，用于改进稿件，\textbf{不判断文本是否由 AI 生成，也不给出任何“AI 生成概率”}。“第 N 行”指 \texttt{main.tex} 源文件行号，“S 第 N 行”指 \texttt{supplementary.tex} 行号，可能与 PDF 行号不同。数值核对只用稿件自身给出的数字。引用文献只核对了 \texttt{references.bib} 中的题名，没有检索全文，所以凡是涉及“引用是否支持主张”的判断都标为“\verify”。有两处判断另外对照了仓库代码（条目 6、条目 41），已在相应条目中注明。建议修改文本中的红色方括号是需要作者补充的信息，本报告不代填任何数字或事实。
\end{abstract}

\tableofcontents
\newpage

%=====================================================================
\section{整体观察}

\subsection{材料与主要写作特点}
稿件提出 M4Fed 框架：各节点用 TabPFN 补全元数据，用 MedMD（文本条件 3D 潜扩散模型）生成合成 MRI，用 SwinMTL（Swin UNETR 多任务预测器）估计脑龄，客户端通过 QLoRA 微调。实验在 OpenMind T1w 数据上进行，包括主比较、组件消融、健康状态消融、超参数与预处理敏感性、跨站点、生成质量、通信量实测和集中式参照。实验部分信息量大，大多数数值可复算。主要写作特点如下。

\begin{enumerate}[leftmargin=*]
\item \textbf{“三类资源 / 三类约束”的表述反复出现（问题类型 2、6）。}“私有数据、公开数据、基础模型”与“隐私、通信、边缘算力”这组表述出现在摘要、第102、104、106、111行、表1、第159、188、200、684、694、703行，句式和措辞高度相似。其中第694行（实践意义段）和第703行（结论）基本是引言的换词复述，没有增加新的信息。这是全文最集中的问题。
\item \textbf{结论范围大于证据（问题类型 3、9）。}摘要、引言第141行、表1和结论写的是无条件的“优于 SOTA”“通信量降低 42.6\%”“边缘参数负担降低 80\% 以上”“同时提升精度、降低通信与算力”。正文中这些结论都带有条件：只有服务器端模型优于基线；42.6\% 只在非默认的 $\rho_f=0.01$ 下成立，默认 $\rho_f=0.1$ 时上传量是 FedAvg 的 5.55 倍；80\% 只统计了域模型的可训练参数，没有计入客户端额外承担的生成模型训练和采样（第692行局限 (iv) 也承认了这一点）。第6.8节和第688行对这些条件写得清楚、坦率，但这些条件没有回传到摘要、表1和结论。
\item \textbf{“结果 → 两条可能原因 → 结论性收束”的固定段落模式（问题类型 2、7）。}第455、517、574、593、604、643、645、686行几乎都按这一模式展开，反复使用“We see two plausible reasons”“A possible / plausible explanation is that”“Another possible explanation”“Interestingly”“It is worth noting”等句式。问题不只是句式单调：有些段落先列出推测，下一句就用“Hence / Therefore / In summary”给出确定结论（第604、645、686行），推测因此被写成了事实。
\item \textbf{一致性问题集中在“默认配置结果”“评估单位”和“域模型输入”三处（问题类型 8、9）。}同一默认配置下的 M4Fed 服务器 MAE 在表5、表7、表8中分别为 8.87、9.22、9.376；训练/测试划分的单位（扫描还是受试者）没有说明；式 (2) 写作 $f_\theta(x,m)$，但 SwinMTL 只以 MRI 为输入。
\item \textbf{文献综述以罗列为主（问题类型 8）。}第106行用“First / Second / Third / Finally”逐条介绍，第163、167行用“On the … side / On the other hand / Besides / More recently / Then”串联文献，对文献之间的比较与综合较少。第159行“None combines …”的断言与表1中 GPT-FL 的标注相矛盾。
\item \textbf{做得较好的部分。}第6.8节的通信量推算全部可复算，并主动指出默认配置下 M4Fed 的上传量更大；第688行把“算力节省无条件成立、通信节省有条件成立”区分得很清楚；第6.9节集中式实验写明了检验方法、样本量和 $p$ 值，并对小样本格子做了提醒；第602行用年龄方差解释 $r$ 和 $R^2$ 下降，论证具体。这几处可以作为全文统一修改的参照。
\end{enumerate}

\subsection{已复算一致的数值（无需修改）}
以下数值用稿件自身给出的数字复算后一致：
$3{,}622+4{,}571-875=7{,}318$；
第448行 $13.7\%$、$1.8\times$；
第455行 $5.4$--$5.6\%$；
第480行 $74\%$、$21\%$；
第515行 MAE 区间 $8.90$--$9.38$，第690行“约 $5\%$”；
第517行 $2.0\%$；
第574行注册实验的全部百分比（SSIM $+29.1\%$/$+16.7\%$，LPIPS $-20.6\%$/$-15.4\%$，FID $-26.7\%$/$-12.1\%$，KID $-10.7\%$，MAE $-4.2\%$）；
第593行 $2.4\%$、$6.8\%$；
第602行 $6.2\%$；
第645行 $-20.8\%$、$-32.4\%$；
第651--653行通信与参数量推算（$4\times9{,}247{,}036\approx37.0$~MB，$96^3\times4=3{,}538{,}944$~B，$484{,}736+1{,}185{,}034=1{,}669{,}770$，$-81.9\%$，$-92.6\%$，$58\beta\approx205.3$~MB，$6\beta\approx21.2$~MB，$5.55\times$，$42.6\%$，平衡点 $\approx10.5$ 卷 / $\rho_f\approx0.018$，1B 参数时 $19.5\times$）；
第677行女性 $307$、男性 $256$、缺失 $8$，合计 $571$，按人数加权的总体 MAE $\approx3.50$，与第675行的 $3.51$ 一致；
表9（第6.5节）“End-to-end”服务器列与表7最后一行、第645行低噪声实验的基线数值（$6.18$、$305.70$、$0.278$）一致。

\subsection{问题分布概览}
\begin{center}
\small
\begin{tabular}{p{0.38\linewidth}p{0.52\linewidth}}
\toprule
问题类型 & 主要集中位置（条目编号） \\
\midrule
1 空泛与套话 & 12、15、35、39 \\
2 机械重复（句式/结构） & 4、5、11、全局模式（第三部分） \\
3 结论放大 & 2、8、9、21、23、26、27、30、35、37、39、40 \\
4 含糊、术语堆砌、逻辑跳跃 & 6、13、14、24、25、29、32、36、40 \\
5 表达不自然 & 3、5、10、23、29、33、34、41 \\
6 重复与信息增量不足 & 4、35、38、39、40 \\
7 事实/推测边界不清 & 14、16、17、24、27、28、29、30、31、36 \\
8 文献与方法细节 & 5、7、11、16、17、18、19、20、33 \\
9 内容一致性 & 1、2、6、7、8、13、16、22、23、24、32、41 \\
\bottomrule
\end{tabular}
\end{center}

%=====================================================================
\section{逐条定位}

%---------------------------------------------------------------------
\subsection{摘要与引言}

\begin{finding}{摘要中的样本量“7318 subjects”与任何单个实验的样本量都不对应}
\loc{摘要，第84行；对照表3（第360行）、第347、482行}
\orig{To train and validate M4Fed, we performed brain-age estimation on 7318 subjects of public datasets.}
\ptype{9 内容一致性；3 结论放大}
\basis{表3中的 $7{,}318$ 是“未筛选队列（$3{,}622$）”和“健康对照队列（$4{,}571$）”去重后的并集，只出现在统计行。正文中的实验分别在其中一个队列上进行（主比较、消融、敏感性用未筛选队列，注册、标准化、低噪声和集中式实验用健康对照队列），没有任何一个实验同时使用 $7{,}318$ 名受试者。另外，数据来源是单一的 OpenMind 语料库，“public datasets”的复数写法也需要说明（OpenMind 本身由多个 OpenNeuro 子数据集汇总而成）。读者会把 $7{,}318$ 理解为单次训练/验证的规模，这会高估实验规模。}
\conf{高}
\normal{报告研究涉及的总受试者数是常见做法，但需要写明这是两个队列的并集。}
\advice{改写}
\revise{We evaluated M4Fed on brain-age estimation using two cohorts drawn from the OpenMind corpus (3,622 and 4,571 subjects; 7,318 unique subjects in total).}
\end{finding}

\begin{finding}{摘要中“surpasses state-of-the-art”“42.6\%”“over 80\%”缺少适用条件}
\loc{摘要第84行；同类表述见第141、703行与表1第133行}
\orig{M4Fed surpasses state-of-the-art approaches while reducing communication payload by 42.6\% and edge-side parameter burden by over 80\%.}
\ptype{3 结论放大；9 内容一致性}
\basis{(1) 表5中只有 M4Fed 服务器端在四项指标上最优；M4Fed 客户端的 MAE（$10.97$）、RMSE（$17.19$）和 $R^2$（$-0.184$）都不如 FedAvg、FedProx 和 GPT-FL，只有 Pearson $r$ 最高。(2) 第653行写明，默认 $\rho_f=0.1$ 时 M4Fed 的单轮上传量是 FedAvg 的 $5.55$ 倍，$42.6\%$ 只在 $\rho_f=0.01$ 时成立，而且只统计了单个客户端的上行流量。(3) “over 80\%”对应的是域模型可训练参数量的降幅（$81.9\%$），没有计入客户端还要训练 MedMD 适配器、运行扩散采样（$\rho_l=1.0$，50 步 DDIM）、存放冻结的 VAE、BiomedCLIP 和 TabPFN；第692行局限 (iv) 也承认生成模型带来了额外的节点训练成本和内存占用。“parameter burden”这一说法本身也比较含糊，读者容易理解为显存或计算量。三处结论合在一起，会让读者以为同一配置下精度、通信、算力三项同时改善。}
\conf{高}
\normal{摘要需要精炼，省略条件属于常见写法；但这里省略的条件会改变结论的方向（默认配置下通信量反而增加），已经超出正常简化的范围。}
\advice{改写}
\revise{The M4Fed server achieves the lowest age-prediction error among all compared methods (MAE 8.87 vs.\ 10.28 years for the best baseline), while M4Fed clients attain the highest Pearson correlation but not the lowest error. QLoRA reduces the clients' trainable domain-model parameters by 81.9\%. Because the per-round upload scales with the synthesis ratio rather than the model size, it is 42.6\% smaller than that of FedAvg at a synthesis ratio of 0.01, and it would be smaller still for larger domain models.}
\end{finding}

\begin{finding}{摘要中的语法和并列结构问题}
\loc{摘要，第84行}
\orig{Consolidating these resources however faces obstacles, including 1) privacy regulations restrict centralization of patient data, 2) large-scale parameters exchanging incurs high communication costs, and 3) limited computing resources … In addition, to exploit imaging and metadata jointly, M4Fed adapts a multi-modal medical generative model, a multi-task predictor, and integrates an imputation model …}
\ptype{5 表达不自然}
\basis{“including”后面接的是三个完整的从句，语法上不成立；“large-scale parameters exchanging”搭配不自然；“adapts A, B, and integrates C”并列结构失衡；“3) employs”与前两项的动词形式也不统一。这些属于语言问题，不涉及实质内容。}
\conf{高}
\normal{属于常见的非母语写作问题，与文本来源无关。}
\advice{改写}
\revise{Consolidating these resources faces three obstacles: (1) privacy regulations restrict the centralization of patient data; (2) exchanging large model parameters incurs high communication costs; and (3) the limited computing resources of individual hospital nodes cannot support full-scale local training. … To exploit imaging and metadata jointly, M4Fed adapts a multi-modal medical generative model and a multi-task predictor, and integrates an imputation model that reconstructs missing metadata.}
\end{finding}

\begin{finding}{“三类资源 / 三类约束”在全文反复以相近句式复述}
\loc{摘要；第102、104、106、111行；表1；第159、188、200、684、694、703行}
\orig{（第111行）… strikes a balance among the three constraints of privacy protection, the communication burden …, and the computational burden on edge nodes, so as to efficiently and collaboratively exploit the three types of resources of private data, public data, and foundation models.\\
（第694行）… usually faces three fragmented training resources: private data …, public datasets …, and large-scale pretrained foundation models. Exploiting all of them together is hindered by three barriers …\\
（第703行）… reconciles three resources previously difficult to consolidate … M4Fed can simultaneously alleviate three practical challenges …}
\ptype{2 机械重复；6 信息增量不足}
\basis{同一组“三资源 + 三约束”在引言中占据了第102、104、106、111四段，在问题定义（第188行）、方法概述（第200行）、讨论（第684、694行）和结论（第703行）中又分别出现，句式基本是“three … namely A, B, and C”。引言内的重复有一定的铺垫作用，但第694行实践意义段几乎完全复述了引言第102--104行，第703行结论的前两句又复述了一遍，后文的复述没有增加证据、边界或新的解释，削弱了讨论和结论的信息量。}
\conf{高}
\normal{摘要、引言和结论各概括一次框架属于正常结构；问题在于复述次数过多，且后文复述没有结合实验结果加以限定。}
\advice{精简}
\revise{建议引言保留一次完整表述（第102--104行）；第111行、第200行改为简短回指（如 “the three constraints introduced above”）；第694行整段删去，或改为结合实验条件的具体陈述（见条目39）；结论只保留结合结果的一句概括（见条目40）。}
\end{finding}

\begin{finding}{引言第106行：逐条罗列式综述、“thereby”连用、搭配不自然，且自造术语写成了既有术语}
\loc{第106行}
\orig{In recent years, a series of methods proposed by researchers have gradually been converging into a technical line that may be collectively termed multi-source federated learning … thereby making it possible … This alleviates the predicament in which edge nodes … and it thereby forges an effective link between foundation models and private data. … Before entering the federated loop, it trains the downstream model … thereby further mobilizing the knowledge embedded within foundation models …}
\ptype{2 机械重复；5 表达不自然；8 文献综述}
\basis{(1) 段落按“First / Second / Third / Finally”逐项展开，每项都是“方法做了什么 + thereby 带来什么好处”，“thereby”一段内出现三次。(2) “forges an effective link”“mobilizing the knowledge”“alleviates the predicament”属于修辞性搭配，含义不具体。(3) “may be collectively termed multi-source federated learning”读起来像是在介绍一个已有术语，但没有给出出处；如果这是作者自己提出的概括，应明确说明，以免读者误以为该术语已有既定定义（是否已有文献使用该术语：\verify）。(4) “some of the latest studies … it trains”主语单复数不一致。(5) 这一段只介绍各方法“做了什么”，没有比较它们各自的局限；局限要到第111行和表1才出现，综述与缺口之间的衔接较弱。}
\conf{高}
\normal{按时间或技术演进顺序综述属于常见写法；问题在于句式过于统一，且缺少比较。}
\advice{改写}
\revise{We use the term \emph{multi-source federated learning} for methods that, beyond the private data of participating clients, also exploit public data or pretrained models. Classical federated learning \citep{…} aggregates model weights instead of raw data but uses only private data. FedDF \citep{…} adds an unlabeled public dataset as a distillation medium. FedPEFT and OpenFedLLM \citep{…} bring in pretrained foundation models and train only low-rank adapters, which lowers client-side computation but still transmits adapter weights every round. GPT-FL \citep{…} additionally uses a generative model to synthesize labeled data for server-side warm-up, yet still transmits the full model in each subsequent round.}
\end{finding}

\begin{finding}{“jointly exploit imaging and metadata”与域模型的实际输入不一致}
\loc{第111行（末句）、第145行（贡献2）、式 (2)（第185行）；对照第288行}
\orig{（第111行）… allowing M4Fed to exploit imaging and metadata jointly rather than treating them as independent modalities.\\
（第185行）$\mathcal{L}_{\mathrm{MT}}(f_\theta(x,m), y)$\\
（第288行）The domain model SwinMTL predicts age and a set of auxiliary clinical fields from MRI volumes.}
\ptype{9 内容一致性；4 概念模糊}
\basis{式 (2) 把域模型写成 $f_\theta(x,m)$，即同时以影像和元数据为输入；第288行和第4.4节的描述则表明 SwinMTL 只以 MRI 为输入，元数据（性别、健康状态等）是预测目标。仓库代码 \texttt{models/swin/swin\_multitask.py} 中 \texttt{forward(self, mri)} 也只接收影像（此项为对照代码所得，代码本身不属于稿件材料）。所以“multi-modal”实际上只体现在 MedMD 的文本条件生成中，元数据通过合成数据间接影响域模型。按现在的写法，读者会以为域模型做的是影像 + 元数据的多模态融合预测；如果 $m$ 中含有年龄，还会让人担心标签泄漏。}
\conf{高}
\normal{不属于表达习惯问题，是符号与方法描述之间的实质不一致。}
\advice{改写}
\revise{式 (2) 改为 $\mathcal{L}_{\mathrm{MT}}(f_\theta(x), y)$。第111行末句改为：“… so that clinical metadata, after imputation, conditions the synthesis of training volumes, while the domain model itself predicts age and auxiliary fields from imaging alone.” 贡献2与标题中的“Multi-Modal”也建议同步说明多模态体现在哪个环节。}
\end{finding}

\begin{finding}{第159行“None combines …”与表1中 GPT-FL 的标注相矛盾}
\loc{第159行；表1第131行；第371行}
\orig{Each of these methods, however, enlarges the resource pool along a single axis only. None combines private client data, public repositories, and pretrained foundation models within one framework …}
\ptype{9 内容一致性；8 文献综述}
\basis{表1中 GPT-FL 在 Private data、Public data、Foundation model 三列均为 \checkmark；第371行也写道“GPT-FL further exploits public data together with the pretrained generative foundation model”。因此“None combines …”和“along a single axis only”与作者自己的表格相矛盾。从表1看，GPT-FL 与 M4Fed 的区别在“约束”一侧（通信 $\triangle$、边缘算力 $\times$），而不在“资源”一侧。}
\conf{高}
\normal{不属于正常表达，是可以直接核对的自相矛盾。}
\advice{改写}
\revise{Among these methods, only GPT-FL draws on all three resources, but it still transmits the full model in every federated round and requires full-backbone training at the clients (Table~\ref{tab:positioning}). No existing method, to our knowledge, exploits all three resources while also addressing the communication and edge-compute constraints.}
\end{finding}

\begin{finding}{表1中 M4Fed 在“通信”“边缘算力”两列标为无条件 \checkmark}
\loc{表1第133行及脚注第137行；对照第653、688、692行}
\orig{M4Fed (ours) \& \checkmark \& \checkmark \& \checkmark \& \checkmark \& \checkmark \& \checkmark}
\ptype{3 结论放大；9 内容一致性}
\basis{表1给 FedPEFT/OpenFedLLM 的通信列标了 $\triangle$，理由是“payload … transmitted every round”；但 M4Fed 同样每轮传输合成数据，而且默认配置下单轮上传量是 FedAvg 的 $5.55$ 倍（第653行）。边缘算力方面，M4Fed 客户端除了 QLoRA 微调域模型，还要 QLoRA 微调 MedMD、为全部私有样本做扩散采样（$\rho_l=1.0$），第692行 (iv) 也承认了额外的训练成本和内存占用。表1的打分标准对基线比对 M4Fed 更严格，评审很容易发现这一点。}
\conf{高}
\normal{定位表通常会突出本方法的优势，但对比列的判断标准应当一致。}
\advice{改写}
\revise{建议在 M4Fed 的通信列和边缘算力列加脚注，例如：“$^{\S}$Per-round payload scales with the synthesis ratio $\rho_f$ rather than with model size; it is smaller than FedAvg's only for $\rho_f \lesssim 0.018$ with the 9.25M-parameter domain model used here (Section~\ref{sec:comm-empirical}). $^{\P}$Clients train only low-rank adapters, but additionally run diffusion sampling (Section~\ref{sec:discussion}, limitation (iv)).” 或者把这两格改为 $\triangle$ 并在脚注中说明。}
\end{finding}

\begin{finding}{第141行和贡献3中“outperforms”“comprehensive”的程度词}
\loc{第141行、第146行}
\orig{The experimental results show that M4Fed outperforms the other baselines … We conduct a comprehensive empirical evaluation …}
\ptype{3 结论放大；1 空泛}
\basis{“outperforms the other baselines”的问题同条目2，只对服务器端成立。“comprehensive”方面：实验只有两个客户端，每个配置看起来只运行了一次（稿件没有说明重复次数或随机种子），主比较没有显著性检验。作为评价性形容词，“comprehensive”没有具体内容支撑。}
\conf{中}
\normal{“comprehensive evaluation”是论文中很常见的写法，单独出现问题不大；与条目2合在一起，会加重结论放大的印象。}
\advice{精简}
\revise{The M4Fed server achieves the lowest prediction error among all compared methods, and M4Fed clients achieve the highest Pearson correlation … We evaluate M4Fed against classical and recent federated learning methods on OpenMind, together with ablation, sensitivity, preprocessing, and cross-site experiments.}
\end{finding}

%---------------------------------------------------------------------
\subsection{相关工作}

\begin{finding}{第156行超长句，以及“now standard”与本文模型规模不一致}
\loc{第156行}
\orig{Both protocols transmit the full parameter set … an assumption that is largely benign for compact convolutional classifiers but becomes a first-order bottleneck once the aggregated model is a volumetric foundation network with tens or hundreds of millions of parameters \citep{ren2025fedfm}, as is now standard in medical image analysis; model compression, gradient sparsification, … have all been proposed as partial remedies …}
\ptype{5 表达不自然；9 一致性}
\basis{这一句约 90 个词，用分号连接了两个独立论点，读起来负担重。另外，“tens or hundreds of millions of parameters … as is now standard”与本文实际使用的域模型（$9.25$M）和生成模型（$7.54$M）不符；正因为模型较小，默认配置下 M4Fed 的通信量才大于 FedAvg（第653行）。作者在这里强调的“瓶颈”规模，本文实验并没有覆盖到。“as is now standard”是概括性判断，引用 \texttt{ren2025fedfm} 是否支持这一点：\verify。}
\conf{中}
\normal{长句在方法类论文中常见；规模上的不一致属于实质问题。}
\advice{改写}
\revise{Both protocols transmit the full parameter set of the model at every round. This is acceptable for compact classifiers but becomes a bottleneck for volumetric foundation models with tens to hundreds of millions of parameters \citep{ren2025fedfm}. Model compression, gradient sparsification, … have been proposed as partial remedies \citep{…}. [建议在本节或第6.8节明确：本文实验所用模型规模较小，大模型情形下的收益来自推算。]}
\end{finding}

\begin{finding}{第163、167行：用连接词串联文献，缺少比较；部分转述超出被引文献的范围}
\loc{第163行、第167行}
\orig{（第163行）… yielding representations that transfer strongly to downstream segmentation and, as we exploit in this work, regression tasks … On the language-vision model side, … On the other hand, tabular foundation model TabPFN … Besides, the MONAI framework …\\
（第167行）Diffusion models … is a paradigm … More recently Kazerouni et al. … Trinh et al. show … Then Mirowski et al. instantiate this direction concretely …}
\ptype{2 机械重复；8 文献综述；5 表达不自然}
\basis{(1) 两段都是“文献 A 做了 X；文献 B 做了 Y；Besides / Then 文献 C …”的罗列，“On the other hand”在这里并不表示对比，用法不当。(2) “More recently Kazerouni et al.（2023）”放在 Rombach 2022 之后、Trinh 2025 和 Mirowski 2026 之前，“More recently / Then”暗示的时间顺序与文献年份不完全对应。(3) “transfer strongly to … regression tasks”：Tang 等人的预训练数据是 CT，本文把它用于 MRI 回归，跨模态迁移的效果需要由本文的消融（w/o pre-training）来支撑，不宜写成既有结论。(4) “Diffusion models … is”主谓不一致。(5) 本节没有说明这些模型之间的关系，也没有说明它们为何适合作为 M4Fed 的组件。}
\conf{中}
\normal{相关工作逐条介绍属于常见写法，但不当的连接词和缺少综合会显得机械。}
\advice{改写}
\revise{（第163行示例）Tang et al.\ \citep{…} pretrain a Swin Transformer encoder on about 5,000 unlabeled CT volumes; whether these representations transfer to MRI-based age regression is examined by our ablation in Section~\ref{sec:ablation-results}. M4Fed uses three further pretrained components, each chosen for a specific role: BiomedCLIP \citep{…} as the text encoder, because …; TabPFN \citep{…} for metadata imputation, because it requires no per-dataset tuning; and the MONAI implementations \citep{…} of the autoencoder, denoising U-Net and Swin UNETR encoder.}
\end{finding}

\begin{finding}{第171行：“collectively affirm … Therefore”的推论较空泛}
\loc{第171行}
\orig{These efforts collectively affirm the effectiveness of parameter-efficient tuning for edge nodes within federated settings. Therefore, M4Fed applies it to a 3D vision backbone and to a diffusion U-Net …}
\ptype{1 空泛与套话；4 逻辑跳跃}
\basis{“collectively affirm the effectiveness”是一般性概括，前文列举的三项工作分别针对 LLM、2D 分类等场景，并不能直接说明 PEFT 对 3D 体数据回归和扩散 U-Net 同样有效；“Therefore”把“别人有效”直接推到“所以我们采用”，跳过了这一差异。更具体的写法是说明本文与这些工作的不同之处。}
\conf{中}
\normal{用“现有工作表明 X 有效，因此我们采用 X”过渡是常见写法，问题不严重。}
\advice{改写}
\revise{These studies apply parameter-efficient tuning to language models and 2D image classifiers. M4Fed extends it to a 3D vision backbone and a 3D diffusion U-Net, using the adapters to …}
\end{finding}

%---------------------------------------------------------------------
\subsection{问题定义与方法}

\begin{finding}{“granularity”没有定义；约束 (b) 的前提与实验所处的规模不符}
\loc{第188行、第200行}
\orig{（第188行）(b) the communication budget per round is limited, making full-model transmission of a foundation-scale $\theta$ infeasible … M4Fed satisfies … (b) by the adjustable synthesis quantity and granularity instead of scaling with $|\theta|$ …\\
（第200行）… synthetic-data aggregation whose volume and granularity are user-configurable.}
\ptype{4 概念模糊；9 一致性}
\basis{(1) “granularity”出现两次，但全文没有定义它指什么（分辨率？每卷体素数？每轮发送的子集？），也没有对应的超参数或实验；正文实际可调的只有 $\rho_f$。(2) 约束 (b) 设定的是“foundation-scale $\theta$ 的全模型传输不可行”，但实验中的域模型只有 $9.25$M 参数（$37.0$~MB），按第653行，FedAvg 在这一规模下的单轮上传量反而小于默认配置的 M4Fed。也就是说，问题定义所设定的约束在实验中并不成立，需要说明大模型情形下的结论来自推算。(3) “M4Fed satisfies (a)”的说法较强：合成数据是否泄露隐私，本文没有评估（见条目14、17）。}
\conf{高（granularity）；中（其余两点）}
\normal{在问题定义中设定约束是正常写法，但需要与实验场景对应。}
\advice{改写}
\revise{… M4Fed addresses (a) by exchanging synthetic samples instead of raw data or weights, (b) by making the per-round payload depend on the synthesis ratio $\rho_f$ rather than on $|\theta|$, and (c) by …. [如果“granularity”确有所指，请明确定义并给出对应参数；否则删去。] [建议补充：“In our experiments the domain model has 9.25M parameters; the advantage under constraint (b) for foundation-scale models is extrapolated in Section~\ref{sec:comm-empirical}.”]}
\end{finding}

\begin{finding}{第202行：“can be exchanged freely”和末句“Therefore …”在逻辑上不成立}
\loc{第202行；对照第692行局限 (v)}
\orig{(iii) Synthetic data $\mathcal{D}^{\mathrm{syn}}$ is generated on-node by the local generative model and can be exchanged freely. Therefore, no incomplete or raw metadata leaves a node in unprocessed form.}
\ptype{7 事实与推测边界；4 逻辑跳跃}
\basis{(1) “can be exchanged freely”把合成数据写成了没有隐私风险的数据，但本文没有做隐私评估（如成员推断、与源样本的相似度检索）；第692行局限 (v) 也承认合成数据可能暴露原始数据的部分特征。前后说法矛盾。(2) “Therefore”前面讲的是“数据分三类”，并不能推出“未处理的元数据不会离开节点”；后一结论实际依赖于第262行（补全在本地进行），以及合成数据所附的文本条件来自补全后的表格这一事实。(3) 同段“On the other hand, data at each node is stratified …”中的“On the other hand”不表示对比，用法不当。}
\conf{高}
\normal{“freely”可能只是想表达“可以在节点间传输”，但字面意思是“不受限制”，在隐私语境下容易引起质疑。}
\advice{改写}
\revise{(iii) Synthetic data $\mathcal{D}^{\mathrm{syn}}$ is generated on-node by the local generative model and is the only data exchanged between nodes. Because imputation is performed locally (Section~\ref{sec:tabpfn}), the metadata strings attached to exchanged volumes are derived from the imputed table rather than from the raw incomplete records. The residual privacy risk of synthetic volumes is discussed in Section~\ref{sec:discussion}.}
\end{finding}

\begin{finding}{第206行“pluggable framework”的可替换性没有经过验证}
\loc{第206行}
\orig{It is worth noting that M4Fed is a pluggable framework rather than a fixed, immutable model. The imputation model, the generative model, and the domain model can each be replaced … so as to accommodate the differing needs of downstream applications and to keep pace with the rapid evolution of foundation models.}
\ptype{3 结论放大；1 空泛}
\basis{实验中三个组件都没有替换过，所以“can each be replaced”是设计意图，而不是已验证的性质；“keep pace with the rapid evolution of foundation models”属于展望性套话。“fixed, immutable”同义重复。}
\conf{中}
\normal{在方法部分说明框架的模块化设计是常见写法，只要措辞上表明这是设计层面的特性即可。}
\advice{精简}
\revise{M4Fed is modular: the imputation, generative, and domain models interact only through synthetic volumes and metadata strings, so each can in principle be replaced by an alternative model. We evaluate only the instantiation described below.}
\end{finding}

\begin{finding}{TabPFN 补全：示例字段“age”与样本筛选条件矛盾，偏差断言和补全效果都缺少证据}
\loc{第258--262行；对照第345行（“Filtering to entries with a non-missing age label”）}
\orig{Consequently, the conditioning distribution becomes biased toward sub-datasets that happen to report that field. … A TabPFN regressor handles continuous fields such as age, weight, and body-mass index …}
\ptype{9 一致性；7 推测写成事实；8 方法细节缺失}
\basis{(1) 第345行说明所有样本都已筛选为年龄非缺失，所以年龄不需要补全；把 age 列为补全字段的例子会让读者担心年龄（预测目标）参与了补全，甚至进入了条件文本。需要说明年龄是否被补全，以及年龄是否作为补全其他字段的输入特征。(2) “Consequently, the conditioning distribution becomes biased …”写成了确定事实，但没有数据或文献支持（属于合理推断，应当这样标注）。(3) 稿件没有报告补全精度（例如在已观测值上做掩码再预测的误差），也没有报告各字段的缺失率；唯一的证据是消融“w/o imputation”，而它的 MAE 与另外两个消融几乎相同（见条目23）。}
\conf{高（第1点）；中（第2、3点）}
\normal{第(2)点是常见的动机论述，加上限定词即可。}
\advice{核实 + 改写}
\revise{… A TabPFN regressor handles continuous fields such as weight and body-mass index \authfill{请核实年龄是否参与补全、是否作为输入特征}, while … Without imputation, a field left blank either propagates an uninformative token or is omitted, which can bias the conditioning distribution toward sub-datasets that report that field. [建议补充各字段缺失率，以及在已观测值上的补全误差。]}
\end{finding}

\begin{finding}{第275行：$\sigma>0.7$ 的理由，以及隐私引用与场景的对应}
\loc{第275行}
\orig{… we use $\sigma > 0.7$ to maximize the diversity of the exchanged synthetic corpus, and thereby the information it carries about the target node's distribution. Such strong noising departs far from the source volume, so it also favors privacy-preserving de-identification \citep{leiblStablePrivacyDiffusionBasedPrivacy2024}.}
\ptype{4 逻辑跳跃；7 证据边界；8 引用对应}
\basis{(1) 噪声越强，合成结果越偏离源样本，保留的源节点信息通常越少；“maximize diversity … and thereby the information it carries about the target node's distribution”把“多样性”和“信息量”画了等号，需要解释为什么更高的多样性反而携带更多关于目标分布的信息，或者改写这一表述。(2) 被引文献 \texttt{leibl…2024} 题名为 “StablePrivacy: Diffusion-Based Privacy Enhancement for Face Image Datasets”，研究对象是人脸图像；它能否支持 3D 脑 MRI 去识别的主张：\verify。本文没有任何隐私度量。(3) 本段长约 330 词，嵌套使用了两套 (i)(ii)… 编号，阅读负担较重。(4) 图4只展示了 $\sigma=0.05$ 和 $0.5$ 的结果，没有展示实际交换时使用的 $\sigma=0.75$。}
\conf{中}
\normal{用文献说明 img2img 噪声强度与隐私的关系是常见做法，但需要说明场景差异。}
\advice{改写 + 核实引用}
\revise{During the federated loop we use $\sigma = 0.75$, so that the exchanged volumes depart substantially from their sources and therefore expose less source-specific detail; a similar argument has been made for diffusion-based de-identification of face images \citep{leibl…}. We do not quantify privacy leakage in this work (Section~\ref{sec:discussion}). [建议把本段拆成“生成流程”和“噪声强度设置”两段。]}
\end{finding}

%---------------------------------------------------------------------
\subsection{实验设置}

\begin{finding}{评估单位（扫描或受试者）、训练轮数与重复次数没有说明}
\loc{第345、365、367、371行；表4；图3、图5图注}
\orig{（第365行）Within each private split, $90\%$ of samples are used for training and $10\%$ for held-out evaluation.\\
（第671行图注）Predicted vs.\ true age of all $571$ test scans.}
\ptype{8 方法细节缺失；9 一致性}
\basis{(1) 第345行写明 $26{,}408$ 次扫描对应 $20{,}107$ 名受试者，也就是说存在同一受试者多次扫描的情况。第365行的“samples”没有说明是按扫描还是按受试者划分；如果按扫描划分，同一受试者的扫描可能同时出现在训练集和测试集中。图注中“571 test scans”与健康对照队列 $4{,}571$ 名受试者的 $10\%$ 量级接近，但第482行又说该队列是按“subject level”划分为四块。需要明确划分单位。(2) 表4只给出了 M4Fed 的 $R_w,R_f,R_l$，没有给出 FedAvg、FedProx、FAN、GPT-FL 的通信轮数、每轮本地 epoch 数，因而读者无法判断各方法的训练预算是否相当。(3) 全文没有说明每个配置运行了几次、是否固定随机种子（第365行的固定种子只针对数据划分）；主表的 $\pm$std 是 2--3 个节点之间的差异，不是重复实验之间的差异。(4) 第371行“a NVIDIA”应为“an NVIDIA”。}
\conf{高}
\normal{不属于表达问题，是复现和公平比较所需的信息。}
\advice{核实 + 补充}
\revise{Sub-datasets are first assigned to roles, and within each private block the split is performed at the \authfill{subject / scan} level, so that \authfill{no subject appears in both training and test sets}. … All baselines are trained for \authfill{rounds} communication rounds with \authfill{epochs} local epochs per round. Each configuration is run \authfill{次数} time(s) with seed(s) \authfill{种子}.}
\end{finding}

\begin{finding}{GPT-FL 复现版本与原方法的差异、FAN 定位需要说明}
\loc{第367、371行；S 第176、181行；对照第106行}
\orig{（第106行）… employ a large language model to expand class labels into prompts, which in turn drive a generative model to synthesize data.\\
（S 第181行）… synthesizes a server-side dataset conditioned on age label seed alone …\\
（第367行）… while GPT-FL's server reuses the generator trained by M4Fed's.}
\ptype{8 方法细节；9 一致性}
\basis{(1) 第106行介绍 GPT-FL 时强调了 LLM 扩写提示词，而补充材料中的复现版本只用年龄标签作为条件，没有 LLM 环节，正文没有说明这一差异。(2) GPT-FL 复用 M4Fed warm-up 阶段训练的生成器，这意味着它也间接使用了 M4Fed 的公开数据预训练；而 FedAvg、FedProx 的域模型训练完全不接触公开数据。因此，M4Fed 服务器相对于这些基线的优势中，有一部分可能来自“多了一份公开数据”，而不完全来自协议本身（见条目21）。(3) 表1把 FAN 列为通信 \checkmark，但补充材料把它描述为“decentralized … perform no server-side aggregation”；FAN 的通信量在第6.8节没有出现，建议补充或在表1中注明。}
\conf{中}
\normal{复现时做简化是常见做法，只要明确说明即可。}
\advice{补充说明}
\revise{Our GPT-FL re-implementation conditions the generator on age-label seeds only and omits the LLM-based prompt expansion of the original method; its generator is the MedMD checkpoint obtained after M4Fed's warm-up, so GPT-FL indirectly benefits from the same public-data pretraining.}
\end{finding}

\begin{finding}{生成质量指标的计算细节缺失}
\loc{第416行；表10}
\orig{… computing PSNR, SSIM, and LPIPS on corresponding axial slices, together with FID and KID computed over the full synthesized-versus-real distribution at each node.}
\ptype{8 方法细节缺失}
\basis{没有说明每个体积取哪些轴向切片、取多少张；FID/KID 使用的特征提取网络（2D Inception 还是医学专用网络）及切片如何输入；评估时使用的噪声强度 $\sigma$；以及每个节点的样本量。表10中的 FID 在 300 左右，读者需要这些信息才能判断数值的可比性。另外，各行指标分别在各自节点的数据上计算，不同行之间的比较还受数据差异影响，这一点也应说明。}
\conf{中}
\normal{属于方法完整性问题，与写作风格无关。}
\advice{补充}
\revise{PSNR, SSIM and LPIPS are computed on \authfill{切片数与位置} axial slices per volume; FID and KID use \authfill{特征网络} features of \authfill{切片/体积数} samples per node, with synthesis performed at $\sigma=$\authfill{数值}.}
\end{finding}

%---------------------------------------------------------------------
\subsection{结果与讨论}

\begin{finding}{主比较：比较单位不对等，没有重复和显著性检验，所有基线 $R^2$ 为负的情况未加讨论}
\loc{第448行；表5}
\orig{M4Fed server is the strongest model by a clear margin on every metric … so the server's accuracy advantage is expected.}
\ptype{3 结论放大；7 证据边界}
\basis{(1) “M4Fed (server)”是单个节点在自己私有块测试集上的结果，基线是 3 个客户端的平均值，两者的测试集不同，比较单位也不同。(2) 服务器额外使用了公开数据和全参数微调（第448行也写了这一点），所以“by a clear margin”有多少来自协议本身、有多少来自额外的数据和算力，从这个表中无法区分；建议补充一个“基线 + 公开数据”或“集中式（仅服务器数据）”的对照，或在文中明确说明这一混杂。(3) 没有重复实验和显著性检验，“by a clear margin”缺乏统计支撑。(4) 除 M4Fed 服务器外，所有方法的 $R^2$ 都为负，说明它们不如直接预测均值；MAE 约 10 年，而集中式参照（第675行）为 3.51 年。这说明当前联邦设定下的整体精度较低，值得在讨论中正面说明，而不是只做相对比较。}
\conf{高}
\normal{“clear margin”在相对差异较大（13.7\%）时是常见说法；但在比较单位不对等且没有统计检验的情况下，应当降低语气。}
\advice{改写 + 补充}
\revise{The M4Fed server attains the lowest error of all methods (MAE 8.87 vs.\ 10.28 years for GPT-FL) and is the only model with a positive $R^2$. This comparison favors the server, which, unlike the baseline clients, also trains on public data with full-parameter updates; the result therefore reflects the combined effect of the protocol and this additional data and compute. All other models have negative $R^2$, i.e., they do not outperform a constant-mean predictor, which indicates that age prediction remains difficult in this federated setting (cf.\ the centralized reference in Section~\ref{sec:centralized-results}).}
\end{finding}

\begin{finding}{默认配置下 M4Fed 服务器的结果在三张表中不一致}
\loc{表5（第442行）、表7（第493行）、表8（第528、533行）；第517、690行}
\orig{表5 M4Fed (server)：MAE $8.87$，RMSE $14.36$，$r$ $0.554$，$R^2$ $0.178$\\
表7 Mixed / Mixed：MAE $9.22$，RMSE $15.11$，$r$ $0.421$，$R^2$ $0.068$\\
表8 $\rho_f=0.1$ 与 lr $=10^{-4}$（即表4默认值）：MAE $9.376$，RMSE $15.172$，$r$ $0.417$，$R^2$ $0.061$}
\ptype{9 内容一致性}
\basis{表4规定的默认值是 $\rho_f=0.1$、MedMD 学习率 $10^{-4}$，第509行说敏感性实验是“swept independently from the default configuration”，因此表8中这两行应当就是默认配置。表7的“Mixed / Mixed”同样是在未筛选队列上用混合数据训练、混合数据测试。但三处结果不同（MAE 8.87 / 9.22 / 9.376，$R^2$ 0.178 / 0.068 / 0.061）。第517行又用 $8.87$ 作为“default partition”的参照。可能的原因包括：不同运行之间的随机波动、不同的队列或划分、不同的代码版本，但稿件没有说明。如果差异来自随机波动，那么 MAE 在 8.87--9.38 之间的波动幅度已经和第515行所称的“robust”区间相当，会影响敏感性分析和主比较的解读。}
\conf{高}
\normal{不属于表达问题。}
\advice{核实}
\revise{请作者核实三处结果是否来自同一配置。若是重复运行，建议报告多次运行的均值 $\pm$ 标准差，并在表8注明默认行；若配置不同，请在表注中说明差异（队列、划分、代码版本等）。在核实之前，本报告不提供替换文本。}
\end{finding}

\begin{finding}{组件消融：“confirming”语气过强，三种消融结果几乎相同，也没有报告服务器端结果}
\loc{第453--455行；表6}
\orig{Table~\ref{tab:ablation} reports client-averaged metrics for each variant as the full-model reference. Every ablation increases MAE relative to the full model, by a fairly uniform $5.4$--$5.6\%$ …, confirming that public-data warm-up, TabPFN-densified conditioning, and pretrained foundation model each contribute positively.}
\ptype{3 结论放大；9 一致性；5 表达}
\basis{(1) 三种性质完全不同的消融（去掉公开数据、去掉补全、随机初始化主干），MAE 分别为 $11.56$、$11.58$、$11.57$，RMSE 为 $18.79$、$18.84$、$18.81$，$R^2$ 的标准差为 $0.001$、$0.002$、$0.000$。这种高度一致与三个组件作用于不同环节的预期不符，可能说明三种变体进入了相同的失效模式（例如预测趋近常数），也可能是记录问题。在核实之前，用它们“confirm”各组件的正向贡献，证据不足。(2) MAE 的差异（约 0.6 年）远小于客户端间的标准差（约 5 年），而且只有两个客户端、没有重复实验，“confirming”语气过强。(3) 消融只报告了客户端结果，而主比较中的优势主要出现在服务器端，消融没有覆盖这一点。(4) 首句“reports client-averaged metrics for each variant as the full-model reference”语义不清；“yields the worse $R^2$”应为“a worse”。(5) 后半段对“w/o public data 时 $r$ 更高”给出了两种推测，属于合理的推测性表述，可以保留。}
\conf{高}
\normal{第(5)点的推测表述正常；第(1)--(3)点属于证据与结论不匹配。}
\advice{核实 + 改写}
\revise{Table~\ref{tab:ablation} reports client-averaged metrics for the full model and each ablated variant. All three ablations increase client MAE, by 5.4--5.6\% (11.56--11.58 vs.\ 10.97 years), although these differences are small relative to the variation across clients. \authfill{请核实三种消融结果几乎一致的原因，并补充服务器端消融结果}.}
\end{finding}

\begin{finding}{健康状态消融：混合测试集的 MAE 高于两个子集，健康测试集的样本量没有报告，论证顺序不清}
\loc{第477--480行；表7}
\orig{This result seems to indicate that when predicting the age of patients with unknown health conditions, mixed data should be used for training, while when predicting the age of healthy individuals, healthy data should be used for training. It is worth noting that the brain age prediction model was used to assess the degree of deviation … Therefore, the main focus should be on the results where the test slice is healthy people. For this part, we flag statistical caution …}
\ptype{9 一致性；4 逻辑跳跃；7 证据边界}
\basis{(1) 混合训练时，混合测试集的 MAE（$9.22$）高于“Ill-only”（$4.45$）和“Healthy-only”（$6.41$）两个子集，这说明混合测试集中还包含大量健康状态缺失的受试者（表3中未筛选队列的 NA 占 $64.4\%$），而且这部分受试者的误差更大。正文没有说明这一点，读者会觉得数字矛盾。(2) 健康子集和患病子集的测试样本量、年龄范围都没有报告，但后文对 $R^2=-39.58$ 的解释（“small healthy-test slice whose age variance is too compressed”）恰恰依赖这些信息；作者手里有这些数据，应当直接报告，而不是推测。(3) 论证顺序是：先给出两条训练建议（“seems to indicate … should be used”），再说明脑龄模型的用途，然后才提示统计上需要谨慎。读者读到建议时还不知道这些数字不可靠。另外，“the brain age prediction model was used to assess”用了过去时，读起来像是在描述本文的做法，但本文并没有计算患者的脑龄差。(4) “It is worth noting”与第206行重复。}
\conf{高}
\normal{第(3)点的结构问题在初稿中常见；第(1)、(2)点属于需要补充的信息。}
\advice{改写 + 补充}
\revise{The mixed test set also contains subjects without a recorded health status (64.4\% of the unselected cohort), which explains why its MAE exceeds that of both labeled subsets. The healthy-only and ill-only test slices contain only \authfill{n} and \authfill{n} subjects with age SD of \authfill{数值} and \authfill{数值} years, so their correlation-based metrics are unreliable: … Within these limits, training on healthy subjects lowers the healthy-test MAE (1.66 vs.\ 6.41 years) but raises the mixed-test MAE (11.15 vs.\ 9.22 years). Because brain-age models are usually trained on healthy subjects and then used to measure deviation in patients, the healthy-test results are the most relevant here.}
\end{finding}

\begin{finding}{第482行：比较的是不同的测试集；用 RMSE 推出“explains more age variance”}
\loc{第482行}
\orig{This value is slightly lower than the mixed-train, healthy-test MAE on the unselected cohort ($6.41$), but the RMSE metric (up to $12.45$) does not show a corresponding gain. Therefore, restricting training to healthy subjects does not necessarily yield a model that explains more age variance on the healthy test slice.}
\ptype{4 逻辑跳跃}
\basis{(1) $6.18$ 和 $6.41$ 来自两个不同队列的不同测试集（健康对照队列的测试集 vs.\ 未筛选队列中的健康子集），两组受试者的年龄分布和样本量都不同，不能直接比较“训练数据”的影响。(2) 健康对照队列按受试者随机均分为四块（包括“server-public”），各块来自相同的站点，这与主实验中按子数据集划分、存在站点异质性的设定不同，作为“reference experiment”时应说明这一差异。(3) RMSE 是误差指标，“explains more age variance”对应的是 $R^2$；用 RMSE 没有改善推出“不能解释更多方差”，指标和结论对不上。(4) “up to $12.45$”措辞含糊，实际上是从 $6.50$ 升至 $12.45$，接近翻倍。(5) “In order to resolve this doubt, we therefore …”中“In order to”和“therefore”语义重复。}
\conf{高}
\normal{不属于表达习惯问题。}
\advice{改写}
\revise{To check whether these effects persist with a larger healthy sample, we train and test M4Fed on the healthy-controls cohort. Note that this cohort is split at random into four subject-level blocks, so all nodes share the same sites, unlike the sub-dataset partition used elsewhere. M4Fed reaches an MAE of 6.18 years, an RMSE of 12.45 years and an $R^2$ of $-0.045$ (Table~\ref{tab:health-ablation}, last row). Because this test set differs from the healthy slice of the unselected cohort, the two results are not directly comparable.}
\end{finding}

\begin{finding}{第515行：“without performance loss”基于单次运行和非单调结果}
\loc{第515行；表8}
\orig{This indicates that M4Fed can maintain performance when significantly reducing the amount of synthetic data used for the federated loop, so that it has potential in terms of largely reducing communication burden without performance loss.}
\ptype{3 结论放大；7 证据边界}
\basis{(1) 每个 $\rho_f$ 只有一次运行，而 MAE 随 $\rho_f$ 的变化不单调（$0.01$: 8.90；$0.1$: 9.38；$1.0$: 8.99），默认值 $0.1$ 反而最差。在随机波动大小未知的情况下（见条目22），不能区分“稳健”与“噪声掩盖了差异”。(2) 另一种同样合理的解释也没有讨论：如果联邦循环中交换的合成数据本身对服务器域模型的贡献很小，那么 $\rho_f$ 自然不影响结果。第686行“synthesis quality and downstream task are only loosely coupled”与这一解释是一致的。这种解释会改变“稳健性”这一结论的含义，值得在文中说明。(3) “has potential in terms of largely reducing … without performance loss”属于冗长表达。}
\conf{中--高}
\normal{“robust to hyperparameters”是敏感性分析的常见结论；但应限定在已测试的范围和运行次数内。}
\advice{改写}
\revise{Within the tested range, server MAE varies between 8.90 and 9.38 years and $R^2$ stays positive, with no monotonic trend in either $\rho_f$ or the learning rate (single run per setting). Lowering $\rho_f$ from 0.1 to 0.01 therefore reduces the federated-loop payload tenfold without a measurable loss in server accuracy in this setting. This insensitivity may also indicate that the synthetic data exchanged in the federated loop contribute little to the server's domain model; we return to this point in Section~\ref{sec:discussion}.}
\end{finding}

\begin{finding}{第517行：“adapt well”与同段数据矛盾，“near-balanced”描述不准确，对 $r$ 符号翻转的解释过于具体}
\loc{第517行；表8 Dirichlet 三行}
\orig{… showing that M4Fed can adapt well to heterogeneous data distributions across nodes. We see two plausible reasons … However, the robustness for non-IID skew also has limits: all three re-partitioned configurations remain somewhat worse than the default partition ($10.64$--$10.86$ vs. $8.87$). … A possible explanation is that a higher degree of non-IID skew makes the data within each node more long-tailed, and M4Fed's local data augmentation reinforces this tendency …}
\ptype{3 结论放大；4 逻辑跳跃；7 推测写成事实}
\basis{(1) 三个重划分配置的 MAE 比默认划分高 $20\%$--$22\%$（10.64--10.86 vs.\ 8.87），$R^2$ 全部为负，“adapt well”与这一事实矛盾，“somewhat worse”也低估了差距。三个 $\alpha$ 之间相差 2\% 只能说明对“偏斜程度”不敏感，不能说明“适应良好”。(2) 对于 3 个节点，对称 $\mathrm{Dir}(1)$ 是单纯形上的均匀分布，抽到的份额可能相当不均衡，称其为“near-balanced”不准确；建议直接报告每个 $\alpha$ 下实际抽到的三个节点份额。(3) 重划分把私有数据合并后重新分配，同时打乱了站点结构，所以“比默认划分差”可能来自站点结构被破坏，而不只是数量偏斜，这一混杂因素需要说明。(4) $\alpha=1.0$ 时 $r=-0.195$（符号翻转），只有一次运行；用“长尾 + 局部增强强化趋势”来解释，属于没有经过检验的具体机制推测，一个更直接的可能是单次运行不稳定。(5) “We see two plausible reasons”与第574行完全相同。}
\conf{高}
\normal{第(4)点用“A possible explanation”加以限定，表述形式本身合规；问题在于解释链条太长，又缺少检验。}
\advice{改写}
\revise{Across $\alpha \in \{0.1, 0.5, 1.0\}$, server MAE varies by only 2.0\% (10.64--10.86 years), so performance is insensitive to the degree of quantity skew. However, all re-partitioned settings are 20--22\% worse than the default sub-dataset partition (8.87 years) and have negative $R^2$; because re-partitioning also mixes sites across nodes, this gap cannot be attributed to quantity skew alone. The drawn node shares were \authfill{各 $\alpha$ 下的实际份额}. Pearson $r$ is negative at $\alpha=1.0$ ($-0.195$); with a single run per setting we cannot tell whether this reflects a systematic effect or run-to-run variability.}
\end{finding}

\begin{finding}{注册实验：解释中引入了表外数据，“support the conclusion of”需要核实}
\loc{第574行}
\orig{… which is consistent with the latent standard deviation moving closer to unit variance after registration ($0.82\to0.89$ before rescaling). … Registration adds little to the clients whose QLoRA updates can be another bottleneck. In summary, affine registration is a low-cost preprocessing step that can further improve generative model and server's domain model. … our results support the conclusion of~\citep{rajabli2025brainage} that some handcrafted preprocessing steps such as registration remain beneficial to the model.}
\ptype{7 证据边界；8 引用对应}
\basis{(1) 潜变量标准差 $0.82\to0.89$ 没有出现在任何表格中，也没有说明在多少样本上估计；作为支持性证据，应当给出来源或放入补充材料。(2) “VAE … pretrained on co-registered BraTS volumes \citep{antonelli2022msd}”：MSD 是数据集论文，可以支持“BraTS 数据经过配准”，但不能支持“该 VAE 是用这些数据预训练的”；后一点应引用所用的 MONAI 预训练权重。补充材料 S 第137行只写了“pretrained for brain MRI generative modeling”，两处描述的具体程度不一致（见条目41）。(3) “QLoRA updates can be another bottleneck”是推测，没有检验。(4) 客户端 MAE 几乎不变（4.970 vs.\ 4.972），服务器 MAE 降低 4.2\%，都只有一次运行，“can further improve … server's domain model”需要加上条件。(5) \texttt{rajabli2025brainage} 题名为 “Brain Age Prediction: Deep Models Need a Hand to Generalize”，主题相符；它是否明确得出了“registration remains beneficial”的结论：\verify。(6) “We see two plausible reasons”与第517行重复。}
\conf{中}
\normal{用两条可能原因解释观察结果是正常写法，第(1)(2)点需要补充依据。}
\advice{补充 + 改写}
\revise{… and the latent standard deviation of the encoder increases from 0.82 to 0.89 after registration (estimated on \authfill{样本数} volumes; Supplementary \authfill{位置}). … In this single-run comparison, registration improves synthesis quality at all nodes and lowers server MAE by 4.2\%, while leaving client MAE unchanged.}
\end{finding}

\begin{finding}{标签标准化：机制解释与 z-score 的数学性质不符，有一处缺动词}
\loc{第593行}
\orig{Figure~\ref{fig:age-zscore-scatter} may show where these gains come mainly from. Without standardization, the predictions of every node mainly around the young-adult majority … Our results suggest that label standardization can partly mitigate this problem because standardization can alleviate the situation where the scale of age labels leads to the model giving more weight to dense samples and neglecting sparse samples.}
\ptype{4 逻辑跳跃；7 推测写成事实；5 表达}
\basis{(1) z-score 是对标签的仿射变换，在 MSE 损失下，它改变的是损失的整体尺度和各节点的均值偏移，并不会改变不同样本之间的相对权重；所以“避免模型偏重密集样本”这一机制与 z-score 的数学性质不符。一个与本文设计直接相关、也更容易检验的解释是：SwinMTL 的年龄头使用 MSE（以年为单位时数值可达数百），而其他头使用 CE，标准化后年龄损失的量级与其他任务更接近，等价于改变了 $\lambda_{\mathrm{age}}$ 和有效学习率。本报告不判断哪种解释正确，只指出现有解释与 z-score 的性质不符，需要作者重新论证或改为推测性表述。(2) “the predictions of every node mainly around the young-adult majority”缺少谓语动词（如“cluster”）。(3) “may show where these gains come mainly from”语气含糊。(4) 把“regression to the mean”等同于“age-dependent bias”在脑龄文献中是常见用法，可以保留。}
\conf{中}
\normal{第(2)(3)点属于语言问题；第(1)点属于实质推理问题。}
\advice{改写}
\revise{Figure~\ref{fig:age-zscore-scatter} shows that, without standardization, the predictions of every node cluster around the young-adult majority, whereas with standardization they spread over a wider range. … One possible reason is that, with the age loss computed in years, the MSE term is much larger than the cross-entropy terms of the auxiliary heads; standardization brings these losses to comparable scales and thus changes the effective weighting of the age task. We did not test this explanation directly.}
\end{finding}

\begin{finding}{跨站点：“conservative”与前一段的论证相矛盾，“generalizes well”和归因都缺少检验}
\loc{第602--604行}
\orig{Hence, the drop of $r$ and $R^2$ mainly reflects the composition of $\mathcal{D}^{\mathrm{cross}}$ rather than a failure of the model … Judged by these point-error metrics, M4Fed generalizes well to unseen sites. … This comparison is conservative, because the baselines are evaluated on test splits drawn from the same sub-datasets as their training data … We attribute this robustness to two parts of the design.}
\ptype{4 逻辑跳跃；3 结论放大；7 推测写成事实}
\basis{(1) 第602行刚论证过，跨站点集合的年龄分布更窄（SD 约 13.3 vs.\ 15.9 年），约 $70\%$ 是 18--24 岁的年轻人；在这种分布上，MAE 和 RMSE 天然更容易偏低。所以用跨站点 MAE 与基线的分布内 MAE 比较，并不是“保守”的比较，测试集差异可能对 M4Fed 有利。“conservative”与作者自己的论证相矛盾。(2) 年龄方差可以解释 $R^2$ 为何转负，但能否完全解释 $r$ 从 0.554 降到 0.202，还需要数据支持；“mainly reflects … rather than a failure of the model”语气过强。(3) “We attribute this robustness to two parts of the design”：没有跨站点的消融（例如 w/o public data 在 $\mathcal{D}^{\mathrm{cross}}$ 上的结果），而所引用的“ablation without public data”是在客户端分布内评估的，且其解释本身也是推测（第455行）。这是在推测的基础上再做推测。(4) 第602行对指标性质的解释写得具体、清楚，是本稿较好的段落，建议保留。}
\conf{高}
\normal{第(3)点改为“may contribute”即属于正常的推测性表述。}
\advice{改写}
\revise{On unseen sites, the server's MAE (9.42 years) and RMSE (13.99 years) remain below the in-distribution values of all baselines. This comparison is, however, not like-for-like: $\mathcal{D}^{\mathrm{cross}}$ has a narrower age distribution, which tends to lower absolute errors, so it shows that point errors do not degrade on unseen sites rather than that M4Fed generalizes better than the baselines. Public-data pretraining at the server and the synthetic volumes exchanged in the federated loop may contribute to this behavior; we did not ablate them on $\mathcal{D}^{\mathrm{cross}}$.}
\end{finding}

\begin{finding}{生成质量：对“服务器不如客户端”的解释与协议描述不符}
\loc{第643行}
\orig{A possible reason is that the server's larger data pool does not actually benefit its generative model. The server warms up MedMD on public data and then broadcasts these warm-up weights unchanged to every client … Hence, compared with the clients, the server gains no extra real-image supervision for MedMD …}
\ptype{4 逻辑跳跃；7 证据边界}
\basis{(1) 按算法1，服务器在 warm-up 阶段用公开数据训练 MedMD，在联邦循环中又用全部客户端的合成数据做全参数微调；客户端从同一个初始权重出发，之后只用服务器合成数据做 QLoRA 微调。这一解释只能说明服务器“没有明显优势”，不能说明服务器为什么“更差”。(2) 各行指标是在各自节点的数据上计算的（第416行“at each node”），节点数据的差异本身就可能造成行间差异，这一点应当先说明。(3) 第643行用“Interestingly”开头，第455行也用了同一个词。}
\conf{中}
\normal{用“A possible reason”限定是正常做法；问题在于解释没有覆盖要解释的现象。}
\advice{改写}
\revise{The server row is slightly weaker than the client average on LPIPS, PSNR and KID. Since each row is evaluated on that node's own data, part of this difference may reflect differences between node datasets. In addition, all nodes start the federated loop from the same warm-up weights, so the server's larger real-data pool is not used to train MedMD beyond the warm-up phase.}
\end{finding}

\begin{finding}{低噪声实验：“CV 领域的噪声区间不能直接迁移到医学影像”的结论没有证据支持}
\loc{第645行}
\orig{Hence, a lower noise strength produces synthetic volumes that are closer to the real data distribution, but this gain in synthesis fidelity does not translate into a better domain model. … Taken together, these results suggest that the range of suitable noise strengths established in the computer vision domain does not transfer directly to medical imaging.}
\ptype{4 逻辑跳跃；9 术语一致性}
\basis{(1) 本段的实际发现是：$\sigma=0.05$ 提高了保真度指标，但下游 MAE 没有改善（6.18 vs.\ 6.21），$r$ 还从 0.249 降到 0.184。也就是说，比 CV 区间（$[0.35,0.70]$）更低的噪声并没有带来下游收益，这与“CV 区间不适用于医学影像”的结论方向相反，至少不支持该结论。实验中也只比较了 $\sigma=0.05$ 和默认设置两种情况，没有扫描区间。(2) 文中“the range … established in the computer vision domain”指的是哪个区间、来自哪篇文献，本段没有给出（第275行引用的是 \texttt{sastry…2024}）。(3) “Lowering $\sigma$ turns the blurry synthesized slices into slices with clear anatomical structure”：$\sigma=0.05$ 时合成结果几乎就是源图像，这一观察是预期之中的，不宜作为方法优点来呈现。(4) “the healthy-subjects-only OpenMind subset”与全文的“healthy-controls cohort”术语不一致。}
\conf{高}
\normal{第(4)点属于术语统一问题。}
\advice{改写}
\revise{… using the healthy-controls cohort. … Thus, a lower noise strength brings the synthetic volumes closer to the real data but does not improve the domain model, possibly because volumes synthesized at $\sigma=0.05$ are nearly copies of their sources and add little diversity. Determining the noise-strength range that best serves downstream training for brain MRI would require a finer sweep of $\sigma$, which we leave to future work.}
\end{finding}

\begin{finding}{通信量实测：只统计了单客户端上行，没有给出总通信量；1B 参数为推算；“Moreover”连用}
\loc{第651--653行}
\orig{Based on the above basic data, the quantitative improvement effects in terms of computing and communication burdens can be analyzed as follows. … For a foundation-scale domain model with $1$B parameters ($4.0$~GB per upload), M4Fed's upload at $\rho_f = 0.1$ would be about $19.5\times$ smaller. Moreover, the sensitivity analysis …}
\ptype{8 方法细节；5 表达}
\basis{(1) 比较只针对“单个客户端每轮上行”，没有包括服务器下行（M4Fed 每轮要广播 $\lceil\rho_f n_s\rceil$ 个合成体积，外加一次性的 30.1~MB 生成模型），也没有按整个训练过程比较总通信量（M4Fed 的 $R_f=3$；基线的轮数在表4中没有给出，见条目18）。摘要中的“reducing communication payload by 42.6\%”实际只对应这一个子项。(2) 1B 参数的情形是推算，前提是合成体积的大小不随模型增大而变化；建议明确标注“extrapolated”。(3) 第一句“Based on the above basic data, the quantitative improvement effects … can be analyzed as follows”是冗长的过渡句，可以删除。(4) “Moreover”在相邻两句中连续出现。本节数值本身全部可复算（见第一部分）。}
\conf{中}
\normal{第(3)(4)点属于语言问题；第(1)(2)点影响摘要结论的范围。}
\advice{补充 + 精简}
\revise{删除首句。末尾改为：“These figures refer to the per-round client upload. Including the server broadcast of $\lceil\rho_f n_s\rceil$ volumes per round and the one-time 30.1~MB generator broadcast, the total traffic over training is \authfill{M4Fed 总量} for M4Fed and \authfill{FedAvg 总量} for FedAvg. For a hypothetical 1B-parameter domain model, the per-round client upload of M4Fed at $\rho_f=0.1$ would be about $19.5\times$ smaller than FedAvg's (extrapolated).”}
\end{finding}

\begin{finding}{集中式参照：“This gap”的指代不明，个别措辞含糊}
\loc{第658、675行}
\orig{The training data and test data are pooled at a single node respectively, and SwinMTL is fine-tuned for $10$ rounds. … the centralized model reaches an MAE of $3.51$, and the predictions are nearly spread along the identity line although there is still bias over age. This gap is expected …}
\ptype{5 表达；9 一致性}
\basis{(1) “This gap”没有明确的比较对象。在同一设定下（健康对照队列、配准、标签标准化），可比的联邦结果应当是第593行的服务器 5.85 年、客户端平均 4.85 年，建议直接写出。(2) 集中式训练中“rounds”的含义不清，应为 epochs；集中式与 M4Fed 的训练总量是否相当，也需要说明。(3) “pooled at a single node respectively”中“respectively”用法不当。(4) “nearly spread along the identity line although there is still bias over age”语义含糊，与第677行对年龄偏差的定量描述重复，可以删去。}
\conf{中}
\normal{属于常见的语言与指代问题。}
\advice{改写}
\revise{All training and test data are pooled at a single node, and SwinMTL is fully fine-tuned for 10 epochs … The centralized model reaches an MAE of 3.51 years, compared with 5.85 years (server) and 4.85 years (client average) for M4Fed under the same preprocessing (Section~\ref{sec:preprocessing-results}). This gap is expected, because …}
\end{finding}

\begin{finding}{讨论第684行：“main finding”只是复述框架设计，“verifies”语气过强}
\loc{第684行}
\orig{Moreover, the component ablation … separately verifies the contributions of public data, pretrained foundation models, and imputed multi-modal data to M4Fed. These results lead to the main finding of this study: M4Fed provides a multi-source federated learning framework for jointly and effectively exploiting the training resources that remain fragmented in the medical domain …}
\ptype{1 空泛；3 结论放大；6 重复}
\basis{(1) “M4Fed provides a framework for exploiting …”描述的是方法的设计目标，而不是实验得出的发现，作为“main finding”是循环表述。(2) “jointly and effectively”中的“effectively”缺乏限定：客户端的点误差没有优于基线，所有客户端模型的 $R^2$ 为负。(3) “separately verifies”的问题同条目23。}
\conf{高}
\normal{讨论开头概括主要结果是正常结构；问题在于概括的是设计而不是结果。}
\advice{改写}
\revise{The main empirical finding is that, in our setting, the M4Fed server achieves a lower age-prediction error than all baselines (13.7\% below GPT-FL), whereas M4Fed clients improve rank correlation but not point error. In the client-level ablation, removing public data, imputation, or pretraining each increases MAE by about 5\%.}
\end{finding}

\begin{finding}{讨论第686行：消融没有去掉 MedMD；“On the contrary”用法不当；推测之后直接用“Therefore”下结论}
\loc{第686行}
\orig{On the generative side, the component ablation … verifies that MedMD can enhance the performance of the domain model by leveraging the imputed multi-modal data. … On the contrary, this might be because the Swin UNETR encoder … can better learn position invariance … Therefore, the impact of registration preprocessing and changing noise intensity is limited, while imputed multi-modal data are more influential.}
\ptype{4 逻辑跳跃；7 推测写成事实}
\basis{(1) 三种消融都没有去掉 MedMD 或合成数据，只去掉了补全；“verifies that MedMD can enhance the performance of the domain model”没有对应的实验支持。要支持这一点，需要“w/o synthetic data”（例如 $\rho_l=0$）的消融。(2) “On the contrary”前后两句并不对立，连接词用法不当。(3) “this might be because …”是推测，下一句却用“Therefore”得出确定结论。(4) “imputed multi-modal data are more influential”是在比较不同实验的结果：补全消融在未筛选队列客户端上的 MAE 变化约 5.5\%，注册实验在健康对照队列服务器上的 MAE 变化为 4.2\%，两者的队列、节点和指标都不同，不能这样比较；而且注册实验中服务器的 $r$ 从 0.249 升到 0.381，变化并不小。(5) 被引文献 \texttt{he2023synthetic}（“Is Synthetic Data from Generative Models Ready for Image Recognition?”）主题相符，具体是否支持“合成数据质量显著影响下游性能”：\verify。}
\conf{高}
\normal{第(5)点属于正常的文献引用，只需核实。}
\advice{改写}
\revise{Our experiments do not isolate the contribution of MedMD's synthetic data to the domain model, because none of the ablations removes the synthetic data. They do show that synthesis quality and downstream accuracy are only loosely coupled: … This contrasts with findings in natural-image recognition that synthetic-data quality substantially affects downstream performance \citep{he2023synthetic}. One possible reason is that the pretrained Swin UNETR encoder, followed by global pooling, is already largely insensitive to spatial misalignment; we did not test this explanation. [建议补充 w/o synthetic data 消融。]}
\end{finding}

\begin{finding}{讨论第690行：预处理结论概括过度}
\loc{第690行末句}
\orig{With respect to preprocessing, the experiment found that registration and label standardization can improve the brain age estimation based on M4Fed.}
\ptype{3 结论放大}
\basis{表9中，注册对客户端 MAE 没有影响（4.970 vs.\ 4.972）；标准化使服务器 $r$ 从 0.381 降到 0.345。两项预处理的效果都因节点和指标而异，而且都只有一次运行，在健康对照队列上进行。笼统写成“can improve”掩盖了这些差异。另外，“the experiment found”用单数，但实际是两组实验。}
\conf{中}
\normal{在讨论中概括为“有助于改善”属于常见写法，但应保留关键差异。}
\advice{改写}
\revise{On the healthy-controls cohort, affine registration mainly improves synthesis quality and server accuracy, whereas label standardization mainly improves client accuracy (single run each).}
\end{finding}

\begin{finding}{局限性部分缺少几项与证据边界直接相关的内容}
\loc{第692行}
\orig{We also note several limitations of the present study. (i) Evaluation is confined to a single modality … (v) We adopt the noise strength … to mitigate the risk that synthetic data expose some features …}
\ptype{6 信息增量不足；7 证据边界}
\basis{现有五条局限写得具体、坦率。但正文中多处影响结论强度的因素没有列入：(1) 每个配置只运行了一次，主比较和敏感性分析没有统计检验；(2) 只有两个客户端，客户端标准差只基于两个值；(3) 服务器同时拥有公开数据和全参数微调，主比较中协议效应与数据/算力效应是混在一起的；(4) 联邦设定下除 M4Fed 服务器外，所有模型的 $R^2$ 为负，与集中式参照（3.51 年）差距较大；(5) 隐私只通过噪声强度间接处理，没有做任何泄露评估。第(ii)条中的“dynamic routing”以及第(iii)条中的公平性激励，与本文结果的关系较远，可以缩短。}
\conf{中}
\normal{局限部分不可能穷尽所有问题；但上述 (1)--(3) 直接影响摘要结论，应当列出。}
\advice{补充}
\revise{(vi) Each configuration was run once with \authfill{种子}, and only two clients were simulated, so the reported differences have not been tested for statistical significance. (vii) The server combines public data with full-parameter fine-tuning; the main comparison therefore does not separate the effect of the protocol from that of this additional data and compute.}
\end{finding}

\begin{finding}{实践意义段（第694行）基本复述引言，“feasible route”超出证据范围}
\loc{第694行}
\orig{Despite these limitations, M4Fed has practical implications for clinical machine learning. In practice, a hospital that wishes to develop a model … usually faces three fragmented training resources … By combining federated consolidation of siloed and public data, cloud-edge collaboration via parameter-efficient fine-tuning, and Synthetic-data-driven aggregation, M4Fed alleviates these three barriers simultaneously. Therefore, M4Fed offers a feasible route for hospitals …}
\ptype{2 机械重复；3 结论放大；6 信息增量不足；1 空泛}
\basis{(1) 这一段的前三句几乎逐项复述了引言第102--104行和第111行的三个设计点，没有引用任何实验结果。(2) “alleviates these three barriers simultaneously”与第653、688行（默认配置下通信量更大）和第692行局限 (iii)(iv)（单机模拟、额外的节点开销）相矛盾。(3) “feasible route for hospitals”：实验是在单机上模拟的，没有在真实医院网络中部署，“feasible”需要限定。(4) 段末又写到“to train machine learning models tailored to their own specific clinical targets”，与段首重复。(5) “Synthetic-data-driven”首字母大写不一致。}
\conf{高}
\normal{讨论中写实践意义是常见做法，但应建立在本文结果和局限的基础上。}
\advice{精简或删除}
\revise{For practitioners, our results suggest the following: when the shared model is small, exchanging synthetic data saves communication only at low synthesis ratios ($\rho_f \lesssim 0.018$ here), whereas the savings grow with model size; client-side QLoRA reduces trainable parameters by more than 80\% but adds the cost of diffusion sampling. Whether these trade-offs hold in a real multi-hospital deployment remains to be tested.}
\end{finding}

\begin{finding}{结论：“simultaneously”和术语堆砌}
\loc{第703行}
\orig{… by combining an imputed multi-modal generative-data-mediated federated protocol with QLoRA-based cloud-edge collaboration. M4Fed can simultaneously alleviate three practical challenges … M4Fed can simultaneously improve predictive accuracy, reduce communication overhead, and lower edge-side compute.}
\ptype{4 术语堆砌；3 结论放大；6 重复}
\basis{(1) “imputed multi-modal generative-data-mediated federated protocol”连续堆叠了五个修饰成分，读者很难解析。(2) “simultaneously”在相邻两句中出现。(3) “simultaneously improve … reduce … lower”需要说明适用条件：在 $\rho_f=0.01$ 时，服务器 MAE（8.90，表8）低于最好的基线，单客户端上行量也比 FedAvg 小 42.6\%，所以“同时”在这一配置下对服务器端有数据支持；但前提是表8的结果与表5可比（见条目22），而且客户端精度并没有改善。现在的写法没有体现这些条件。(4) 前两句又一次复述了“三资源 / 三挑战”。}
\conf{高}
\normal{结论需要概括，但应当与第6.8节的条件保持一致。}
\advice{改写}
\revise{We presented M4Fed, a federated learning framework that combines private client data, public data, and pretrained medical foundation models by exchanging metadata-conditioned synthetic MRI volumes and restricting client training to QLoRA adapters. On brain-age estimation with OpenMind, the M4Fed server achieved lower error than FedAvg, FedProx, FAN, and GPT-FL; at a synthesis ratio of 0.01 it did so with a per-round client upload 42.6\% smaller than that of FedAvg, and QLoRA reduced the clients' trainable domain-model parameters by 81.9\%. M4Fed clients improved rank correlation but not point error. Future work includes …}
\end{finding}

%---------------------------------------------------------------------
\subsection{补充材料}

\begin{finding}{补充材料与正文在 VAE 来源、术语上不一致，另有个别语法问题}
\loc{S 第137行、S 第147行、S 第58行；对照正文第574行}
\orig{（S137）… instantiated from a MONAI AutoencoderKL pretrained for brain MRI generative modeling …\\
（S147）… which is the adaptation needed for federated personalization.\\
（S58）… at client nodes replacing full-model.}
\ptype{9 一致性；5 表达}
\basis{(1) 正文第574行说 VAE 是在“co-registered BraTS volumes”上预训练的，补充材料只说“pretrained for brain MRI generative modeling”。仓库配置 \texttt{configs/base\_config.py} 中的 \texttt{vae\_repo = "MONAI/brats\_mri\_generative\_diffusion"} 与正文说法一致（此项为对照代码所得）。建议在补充材料中写明具体的权重来源，并在正文给出相应引用。(2) “federated personalization”在正文中从未出现，正文的目标是训练各节点自己的 SwinMTL，是否属于“personalization”需要统一表述。(3) S 第58行“replacing full-model”语法不完整。(4) S 第137行“renders 3D diffusion tractable on hospital-grade accelerators”没有给出硬件或显存数据，与正文第167行的说法相同，属于没有量化的断言。}
\conf{中}
\normal{第(3)点属于语言问题。}
\advice{改写}
\revise{（S137）… instantiated from the MONAI \texttt{brats\_mri\_generative\_diffusion} AutoencoderKL checkpoint \authfill{引用}, which was pretrained on BraTS volumes … （S58）… at client nodes, instead of fine-tuning the full model.}
\end{finding}

%=====================================================================
\section{综合修改}

\subsection{优先级最高、最需要修改的段落}
\begin{enumerate}[leftmargin=*]
\item \textbf{默认配置结果不一致（条目22）与消融结果高度一致（条目23）。}这两处需要先回到实验记录核实，它们决定了主比较、消融和敏感性分析的可信度，其他文字修改都应在核实之后进行。
\item \textbf{摘要、表1、引言第141行、结论第703行（条目1、2、8、9、40）。}把正文第6.8节和第688行已经写清楚的条件（只有服务器端最优；42.6\% 对应 $\rho_f=0.01$ 和单客户端上行；80\% 对应域模型可训练参数）同步到这些位置。
\item \textbf{域模型输入与“多模态”的表述（条目6）。}修改式 (2)，并在标题、贡献2和第111行统一说明多模态体现在哪个环节。
\item \textbf{第159行与表1的矛盾（条目7）。}
\item \textbf{讨论第684、686、694行（条目35、36、39）。}这三段集中了“复述设计”“推测之后下结论”“跨实验比较”的问题，建议重写为基于结果的讨论。
\item \textbf{跨站点“conservative”（条目30）、低噪声实验结论（条目32）、Dirichlet“adapt well”（条目27）。}这三处的结论与同段数据相矛盾，修改成本低，收益大。
\item \textbf{评估单位、训练轮数、重复次数（条目18）。}这些是评审一定会问的复现信息。
\end{enumerate}

\subsection{反复出现的表达问题与统一修改原则}
\begin{center}
\small
\begin{longtable}{p{0.30\linewidth}p{0.25\linewidth}p{0.35\linewidth}}
\toprule
反复出现的表达 & 出现位置（行） & 统一修改原则 \\
\midrule
\endhead
“three resources / three constraints / three barriers” & 84, 102, 104, 106, 111, 159, 188, 200, 684, 694, 703 & 引言完整表述一次，其余位置回指或结合结果给出限定；删去第694行的复述。\\
“We see two plausible reasons” / “A possible (plausible) explanation” / “Another possible explanation” & 455, 517, 574, 593, 643, 645 & 只保留有依据的解释；没有检验的解释写一句即可，并注明“we did not test this”；不要在推测之后用“Hence / Therefore”下结论。\\
“It is worth noting that” & 206, 480 & 删去，直接陈述。\\
“Interestingly” & 455, 643 & 删去，或改为说明为什么出乎意料。\\
“On the other hand”（无对比关系） & 163, 202, 448 & 按实际逻辑改为“In addition”或删去；第448行属于正当的对比，可以保留。\\
“thereby” 连用 & 106（3次）、275、347 & 每段最多保留一次，或拆成两句。\\
“Moreover” 相邻连用 & 653 & 删去其中一个。\\
无条件的程度词：“clear margin”“comprehensive”“adapt well”“generalizes well”“effectively”“feasible”“simultaneously” & 141, 146, 448, 517, 604, 684, 694, 703 & 只要正文存在条件，就在同一句里写出条件（配置、节点、指标、运行次数）。\\
“verifies / confirms”用于单次运行、小差异的结果 & 455, 684, 686 & 改为“is consistent with”“suggests”，或直接报告数值。\\
修饰性搭配：“forges an effective link”“mobilizing the knowledge”“clinically grounded”“keep pace with the rapid evolution” & 106, 111, 206 & 改为具体的动作或删除。\\
术语不统一：“healthy-subjects-only OpenMind subset” vs.\ “healthy-controls cohort”；“Synthetic-data-driven”大小写；“federated personalization” & 645, 694, S147 & 全文统一为同一术语。\\
\bottomrule
\end{longtable}
\end{center}

统一原则概括如下：
\begin{enumerate}[leftmargin=*]
\item \textbf{结论与条件写在同一句里。}凡是正文中有适用条件的结论（节点、配置、指标、队列），在摘要、讨论和结论中复述时都要带上这些条件。
\item \textbf{区分三类陈述。}本文直接观察到的结果，用数值陈述；文献支持的结论，加引用；作者的推测，用“may / possibly”，并说明没有检验。不要让第三类陈述后面紧跟“Therefore”。
\item \textbf{不跨实验比较效应大小。}队列、节点或指标不同的实验，不直接比较“哪个因素影响更大”。
\item \textbf{减少段落模板。}结果段不必都采用“结果 → 两条原因 → 总结”的结构；有些段落只需报告结果并指出边界即可。
\item \textbf{不为降低“模板感”而加入口语或不规则表达}，修改的方向是让每句话承载具体的信息。
\end{enumerate}

\subsection{需要补充数据、文献，或需要与上下游核实的判断}
\begin{enumerate}[leftmargin=*]
\item 表5、表7、表8中默认配置结果（8.87 / 9.22 / 9.376）的来源与差异原因（条目22）。
\item 三种消融结果几乎相同、$R^2$ 标准差接近 0 的原因；服务器端消融结果（条目23）。
\item 训练/测试划分单位（扫描或受试者），以及是否存在同一受试者跨集合的情况（条目18）。
\item 基线的通信轮数和本地 epoch 数；每个配置的运行次数和随机种子（条目18）。
\item 健康子集、患病子集、健康状态缺失子集的测试样本量和年龄标准差（条目24）。
\item 各 $\alpha$ 下实际抽到的节点数据份额（条目27）。
\item 年龄是否参与 TabPFN 补全，或是否作为补全其他字段的输入；各字段缺失率和补全误差（条目16）。
\item 潜变量标准差 $0.82\to0.89$ 的估计方式与来源（条目28）。
\item 生成质量指标的切片选择、特征网络和评估时的 $\sigma$（条目20）。
\item 包括服务器下行和全部轮次在内的总通信量（条目33）。
\item 是否补充“w/o synthetic data”消融，以支持第686行关于 MedMD 贡献的说法（条目36）。
\item 以下引用与主张的对应关系需要检索全文核实（本报告只核对了题名）：
  \bibkey{leiblStablePrivacyDiffusionBasedPrivacy2024}（人脸图像去识别，用于支持 MRI 隐私，条目17）；
  \texttt{antonelli2022msd}（数据集论文，用于支持 VAE 预训练，条目28）；
  \texttt{rajabli2025brainage}（是否明确支持“配准有益”，条目28）；
  \texttt{he2023synthetic}（条目36）；
  \texttt{ren2025fedfm}（是否支持“tens or hundreds of millions … now standard”，条目10）；
  \bibkey{sastryDiffAugDiffuseandDenoiseAugmentation2024}（$[0.35,0.70]$ 区间的出处，条目32）；
  \texttt{kocak2026undertrained}（放射组学样本量论文，用于支持“单中心样本过小”这一一般性主张）；
  \texttt{phuntsho2025adaptation}（是否支持“clinically trustworthy predictions”的说法）；
  以及“multi-source federated learning”是否为已有术语（条目5）。全部标为 \verify。
\end{enumerate}

\subsection{说明}
本报告指出的写作特征，包括句式重复、段落模板化、修饰性搭配、结论放大和推测与事实的边界不清，在人工写作中同样常见，\textbf{不能用来证明文本由 AI 生成}，本报告也不对此作任何判断。报告中的建议只针对可读性、论证严谨性和文内一致性。建议修改文本只调整表述，没有新增数据、参数或结论；红色方括号中的内容需要作者根据实验记录补充。
\end{document}
